\section{Comparison with Alternative Approaches}


\subsection{Fine-Tuning}

\textbf{Advantages of Fine-Tuning}:
\begin{itemize}
    \item Optimal performance on single domain (if sufficient data)
    \item No inference-time overhead (parameters encode all knowledge)
\end{itemize}

\textbf{Disadvantages}:
\begin{itemize}
    \item Cost: 667× more expensive (\$12K vs. \$18)
    \item Catastrophic forgetting across domains (Table~\ref{tab:transfer})
    \item Opacity: No interpretability of parameter changes
    \item Non-composable: Cannot combine multiple fine-tunes
\end{itemize}

\subsection{RLHF and GRPO}

\textbf{Vanilla GRPO} improves upon RLHF by eliminating value networks, reducing training instability. However:

\begin{itemize}
    \item Still updates parameters ($\theta \neq \theta_0$)
    \item Advantages are numerical, not semantic
    \item No audit trail for why performance changed
\end{itemize}

\textbf{\DSO{}} extends GRPO into the semantic realm: advantages become human-readable experiences, enabling decentralized curation.

\subsection{Prompt Engineering}

Manual prompt crafting is effective but:

\begin{itemize}
    \item Requires expert knowledge (not scalable)
    \item Static (does not improve with model usage)
    \item No mechanism for community contributions
\end{itemize}

\textbf{\DSO{}} automates prompt refinement via introspection and scales via decentralized governance.

\subsection{RAG}

RAG retrieves factual documents, while \DSO{} retrieves \textit{strategic insights}. Key differences:

\begin{table}[t]
\centering
\caption{RAG vs. \DSO{}}
\label{tab:rag-vs-dso}
\begin{tabular}{lcc}
\toprule
\textbf{Aspect} & \textbf{RAG} & \textbf{\DSO{}} \\
\midrule
Retrieval target & Facts & Strategies \\
Source & External corpus & Model introspection \\
Update mechanism & Manual corpus edit & Automated extraction \\
Governance & Centralized & Decentralized \DAO{} \\
Performance gain & +2--5\% & +10--15\% \\
\bottomrule
\end{tabular}
\end{table}

